\subsection{Estructuras de Control}

\subsubsection{Regla: llaves obligatorias en todo bloque}

\textbf{Regla:} Es obligatorio el uso de llaves \texttt{\{\}} en toda estructura de control, incluso cuando el bloque contenga una única instrucción.

\textbf{Descripción:} La regla de estilo de código \texttt{IDE0011} de .NET establece \texttt{true} como valor predeterminado de la opción \texttt{csharp\_prefer\_braces}, es decir, preferir llaves incluso para una sola línea de código (Microsoft, 2023).

\textbf{Código correcto:}
\begin{lstlisting}[style=csharpstyle]
if (diceRoll == MaximumFaceValue)
{
    GrantExtraTurn();
}
\end{lstlisting}

\textbf{Código incorrecto:}
\begin{lstlisting}[style=csharpstyle]
if (diceRoll == MaximumFaceValue)
    GrantExtraTurn();
\end{lstlisting}

\subsubsection{Regla: rama \texttt{default} obligatoria en \texttt{switch}}

\textbf{Regla:} Toda declaración \texttt{switch} deberá incluir siempre una rama \texttt{default} explícita que gestione los valores no contemplados.

\textbf{Descripción:} McConnell indica que las estructuras de selección múltiple deben contemplar de forma explícita los casos no previstos para evitar que el programa continúe en un estado inconsistente (McConnell, 2004).

\textbf{Código correcto:}
\begin{lstlisting}[style=csharpstyle]
public string DescribeTurnState(TurnState state)
{
    switch (state)
    {
        case TurnState.WaitingForRoll:
            return "Waiting for roll";
        case TurnState.Moving:
            return "Moving";
        case TurnState.Finished:
            return "Finished";
        default:
            throw new ArgumentOutOfRangeException(nameof(state), state, "Unsupported turn state.");
    }
}
\end{lstlisting}

\textbf{Código incorrecto:}
\begin{lstlisting}[style=csharpstyle]
public string DescribeTurnState(TurnState state)
{
    switch (state)
    {
        case TurnState.WaitingForRoll:
            return "Waiting for roll";
        case TurnState.Moving:
            return "Moving";
        case TurnState.Finished:
            return "Finished";
    }

    return string.Empty;
}
\end{lstlisting}
